\section*{Chapitre \ref{ch:qm:brownian-motion} --- Mouvement brownien}

\begin{solution}{exo:qm:brownian-motion:1}
$\Cov(W_2, W_5) = \min(2,5) = 2$~; $\Var(W_5 - W_2) = 3$~; corrélation $2/\sqrt{2 \times 5} = 0.632$.
\end{solution}

\begin{solution}{exo:qm:brownian-motion:2}
Par le principe de réflexion, $\P(M_1 \ge 1) = 2(1 - \Phi(1)) = 31.7\%$.
\end{solution}

\begin{solution}{exo:qm:brownian-motion:3}
$0.20/\sqrt{252} = 1.26\%$ par jour~; $\sqrt5$ fois cette valeur, soit 2.82~\%, par semaine.
\end{solution}

\begin{solution}{exo:qm:brownian-motion:4}
$b = \ln 0.95 = -0.0513$, $T = 0.25$~: en continu $2\Phi(b/0.15) = 73.2\%$~; aux clôtures, le
niveau décalé de $0.0110$ donne 67.8~\%.
\end{solution}

\begin{solution}{exo:qm:brownian-motion:5}
$\E_s[W_t^2] = W_s^2 + (t - s)$, donc $W_t^2 - t$ est une \omterm{def:qm:probability-at-speed:martingale}{martingale}. En notant $\tau$ le temps de sortie,
$\E[W_{\tau\wedge n}^2] = \E[\tau \wedge n]$~; en faisant tendre $n \to \infty$ (convergence dominée à
gauche, monotonie à droite) et en utilisant $\P(W_\tau = b) = a/(a+b)$~: $\E[\tau] = b^2\frac{a}{a+b} +
a^2\frac{b}{a+b} = ab$.
\end{solution}

\begin{solution}{exo:qm:brownian-motion:6}
$\mathcal N(0, \tfrac14)$, à partir de la covariance $s(1 - s)$~; par
la \cref{prop:qm:brownian-motion:between}, $\P(\max \ge 0.5) = \exp(-2 \times 0.5 \times 0.5) = 60.7\%$.
\end{solution}

\begin{solution}{exo:qm:brownian-motion:7}
$\Var(\sum\Delta_k^2) = \sum 2(1/n)^2 = 2/n$~: pour $n = 4\,096$, l'écart type vaut 0.022, et
2\,000 trajectoires simulées donnent 0.022. La valeur de la figure pour 4\,096 points, 1.031, se situe
à 1.4 écart type de 1~: cohérent.
\end{solution}

\begin{solution}{exo:qm:brownian-motion:8}
Contrôler les clôtures quotidiennes revient à évaluer un contrat surveillé chaque jour, et non en continu~: il manque chaque
franchissement entre deux clôtures et sous-estime la probabilité de désactivation de plusieurs points (72~\%
contre 82~\% pour le stop à 2~\% du chapitre). Davantage de trajectoires ne réduisent que l'erreur statistique,
mais pas ce biais, qui décroît comme $\sqrt{\Delta t}$. Ajouter à chaque intervalle la probabilité de franchissement
du \omterm{def:qm:brownian-motion:bridge}{pont brownien}, ou évaluer le contrat continu avec le niveau décalé de $\beta\sigma\sqrt{\Delta t}$
pour vérifier.
\end{solution}

\begin{solution}{pb:qm:brownian-motion:1}
\textbf{1.} $b = \ln 0.98 = -0.0202$~; $\sigma\sqrt T = 0.30\sqrt{21/252} = 0.0866$.
\textbf{2.} $\Phi(-0.0202/0.0866) = \Phi(-0.233) = 40.8\%$.
\textbf{3.} $2 \times 40.8\% = 81.6\%$.
\textbf{4.} Par le principe de réflexion, chaque trajectoire qui touche le stop et termine au-dessus de lui est
associée à son image symétrique, qui termine en dessous.
\textbf{5.} 82.4~\%~: la dérive négative rend l'atteinte légèrement plus probable.
\textbf{6.} 72.0~\% (erreur type de 0.1 point).
\textbf{7.} $0.5826 \times 0.30/\sqrt{252} = 0.0110$ en log-prix~: le stop se comporte comme un stop de 3.1~\%
surveillé en continu.
\textbf{8.} 71.9~\%.
\textbf{9.} 79.6~\% par simulation, 79.6~\% par décalage.
\textbf{10.} Le dépassement, et donc l'erreur, est d'ordre $\sigma\sqrt{\Delta t}$~: diviser
$\Delta t$ par quatre la réduit de moitié (77.1~\%, 79.1~\%, 80.5~\% pour 4, 16 et 64 observations par jour).
\textbf{11.} $\exp(-2 \times 0.01 \times 0.01/(0.09/252)) = 57.1\%$.
\textbf{12.} 81.6~\% (erreur type de 0.08 point), égal à la réponse en continu.
\textbf{13.} Connaissant les clôtures, les franchissements de jours différents sont des bridges indépendants, donc $1
- \prod(1 - p_k)$ est la probabilité conditionnelle d'atteinte~; sa moyenne est la probabilité inconditionnelle.
\textbf{14.} Le bridge~: il est exact, au prix de pas quotidiens, alors que 64 pas par jour
laissent encore un biais de 1 point, pour un coût 64 fois supérieur.
\textbf{15.} 64.1~\%.
\textbf{16.} 89.3~\% en continu, 83.5~\% aux clôtures.
\textbf{17.} Le chiffre en continu, 82~\%~; le modèle ignore les sauts de cours (l'ordre s'exécute sous le stop
lorsque le cours franchit le stop d'un bond), la volatilité non constante et les profils intrajournaliers.
\textbf{18.} Le chiffre des clôtures quotidiennes, 72~\%~: le contrat est surveillé aux clôtures.
\textbf{19.} Résultat nommé~: \emph{le stop-loss}~: un stop de 2~\% sur une action dont la volatilité est de 30~\% est atteint
dans un mois avec une probabilité de 81.6~\% s'il est surveillé en continu et de 72.0~\% s'il est surveillé aux
clôtures~; le décalage de Broadie--Glasserman--Kou donne 71.9~\% et la simulation quotidienne corrigée par le
bridge donne 81.6~\%.
\textbf{20.} Parce qu'une trajectoire brownienne franchit et refranchit un niveau entre deux observations, et
les clôtures ne voient que les franchissements qui subsistent jusqu'à la clôture.
\end{solution}

\begin{solution}{iq:qm:brownian-motion:1}
$W$ est une \omterm{def:qm:probability-at-speed:martingale}{martingale} bornée jusqu'à la sortie, donc $\E[W_\tau] = 0$~: $-1 \times p + 2(1 - p) = 0$,
$p = 2/3$.
\iqlookfor{arrêt optionnel en une ligne.}
\end{solution}

\begin{solution}{iq:qm:brownian-motion:2}
Les accroissements au carré sur une subdivision ont une moyenne égale au temps écoulé et une variance
qui s'annule avec le pas de la subdivision. C'est pourquoi la variance d'un portefeuille couvert s'accumule avec le temps
quelle que soit la fréquence de rééquilibrage, pourquoi la volatilité est mesurable sur une seule trajectoire, et pourquoi
la formule d'Itô comporte un terme du second ordre.
\iqlookfor{le calcul dans $L^2$ et ses conséquences.}
\end{solution}

\begin{solution}{iq:qm:brownian-motion:3}
Par le principe de réflexion, $M_1$ suit la loi de $|W_1|$, donc $\E[M_1] = \sqrt{2/\pi} = 0.798$.
\iqlookfor{réflexion, puis moyenne de la loi demi-normale.}
\end{solution}

\begin{solution}{iq:qm:brownian-motion:4}
Simuler sur la grille dont on a de toute façon besoin et, dans chaque intervalle, multiplier par la probabilité
de non-franchissement du bridge, $1 - \exp(-2(x_k - b)(x_{k+1} - b)/\sigma^2\Delta t)$ (ou tirer le
franchissement avec cette probabilité)~; pour un contrôle rapide, décaler le niveau de $0.5826\,\sigma\sqrt{\Delta t}$.
\iqlookfor{la probabilité de franchissement du \omterm{def:qm:brownian-motion:bridge}{pont brownien}.}
\end{solution}

\begin{solution}{iq:qm:brownian-motion:5}
La probabilité de terminer le mois sous le stop, 41~\%. Un stop est déclenché par le minimum,
et non par la valeur terminale~; par réflexion, le bon chiffre est le double de cette valeur, soit 82~\% en continu, 72~\%
si seules les clôtures comptent.
\iqlookfor{probabilité terminale contre probabilité dépendant de la trajectoire.}
\end{solution}

\begin{solution}{iq:qm:brownian-motion:6}
$(W_1, W_2 - W_1)$ sont des normales centrées réduites indépendantes $(X, Y)$, et l'événement est $X > 0$, $X + Y >
0$~: un secteur d'angle $\pi/2 + \pi/4$ dans le plan, donc la probabilité vaut $(3\pi/4)/(2\pi) = 3/8$.
\iqlookfor{symétrie de rotation de la loi gaussienne.}
\end{solution}
